\documentclass[twocolumn, a4paper, polish, 10pt]{article}
\usepackage[polish]{babel}
\usepackage[utf8]{inputenc}
\usepackage[T1]{fontenc}
\usepackage{secdot}
\usepackage{indentfirst}
\usepackage{graphicx}
\usepackage{color}
\usepackage{sectsty}%pozwala na zmianê typu czcionki paragrafu
\usepackage[usenames,dvipsnames,svgnames,table]{xcolor}
\definecolor{fiolet}{cmyk}{0.47,1,0.07,0.33}
\usepackage[a4paper, hmargin=2cm, tmargin=1.5cm, bmargin=3cm, footskip=2cm]{geometry}
\usepackage{fancyhdr}

% Additional imports
\usepackage[colorinlistoftodos]{todonotes}
\usepackage{amsmath}
\usepackage{amssymb}
\usepackage{physics}
\usepackage{enumitem}
\usepackage{pgfplots}
\usepackage{hyperref}
\usepackage{xcolor}
\usepackage{xurl}
\usepackage{csquotes}

\setlength {\marginparwidth }{2cm}
\pgfplotsset{compat=1.18}
\hypersetup{
    colorlinks=true,
    linkcolor=black,
    urlcolor=blue,
    citecolor=blue
}

\DeclareMathOperator{\EX}{\mathbb{E}}

\pagestyle{fancy}
\lhead{}
\chead{}
\rhead{}
\cfoot{}
\rfoot{\emph{Paweł~Konopka}}
\renewcommand{\headrulewidth}{0pt}
\renewcommand{\thefootnote}{\fnsymbol{footnote}}
\cfoot{\thepage{}}
%\lfoot{\footnotesize{\textit{pawel.konopka@student.uj.edu.pl}}}

\begin{document}

\sectionfont{\centering\large}%kolor czcionki oraz centrowanie tytułu paragrafu
\twocolumn[
\begin{@twocolumnfalse}
\begin{center}
\textbf{\large{Zadanie zaliczeniowe z kursu SMOP I}}
\bigskip
\\Paweł~Konopka\footnotemark \\\textit{Statystyczne metody opracowania pomiarów I}\\
\small{data rozpoczęcia ćwiczenia: 23 grudnia 2025}\\
\end{center}

%Abstrakt
\begin{quote}\small{
\begin{quote}
\begin{center}
\textbf{Streszczenie}
\end{center}
To sprawozdanie jest zadaniem zaliczeniowym kursu \textit{Statystyczne metody opracowania pomiarów I}, bazującym na ćwiczeniu \textbf{M8} z \textit{I pracowni fizycznej} \cite{pracownia_fizyczna_m8}. W ramach zadania otrzymałem \textbf{Zestaw~danych~33}.
\end{quote}
}
\end{quote}
 \bigskip
    \bigskip
\end{@twocolumnfalse}
]
\footnotetext{pawel.konopka@student.uj.edu.pl}

\section{Wstęp teoretyczny}
\subsection{Metoda I}
Pierwszą z metod obliczenia momentu bezwładności tarczy jest pomiar jego masy $M_{\mathrm{w}}$ oraz promienia $R_{\mathrm{w}}$ podstawy, a następnie skorzystanie ze wzoru na moment bezwładności walca \cite{eFizyka_bryla_sztywna}:
\begin{equation}
    J_{\mathrm{w}} = \frac{1}{2} M_{\mathrm{w}} R_{\mathrm{w}}^2
\label{eq:moment_bezwl_walca_promien}
\end{equation}

W przypadku pomiaru średnicy podstawy walca $D_{\mathrm{w}}$ zamiast średnicy ($R_{\mathrm{w}} = \frac{1}{2} D_{\mathrm{w}}$), wzór~(\ref{eq:moment_bezwl_walca_promien}) przyjmuje postać:
\begin{equation}
    J_{\mathrm{w}} = \frac{1}{2} M_{\mathrm{w}} (\frac{1}{2} D_{\mathrm{w}})^2 = \frac{1}{8} M_{\mathrm{w}} D_{\mathrm{w}} ^2
\label{eq:moment_bezwl_walca_srednica}
\end{equation}

\subsection{Metoda II}
Drugą z metod obliczenia momentu bezwładności krążka jest skorzystanie z następującej zależności dla żyroskopu \cite{pracownia_fizyczna_m8}:
\begin{equation}
    mgr = \frac{ J 4 \pi^2 }{T T_{\mathrm{p}}}
\label{eq:zyroskop}
\end{equation}

gdzie:
\begin{itemize}[nosep]
  \item $m$ - dodatkowe obciążenie żyroskopu
  \item $g$ - przyspieszenie ziemskie
  \item $r$ - ramię działania dodatkowego obciążenia
  \item $J$ - szukany moment bezwładności bryły
  \item $T$ - okres obrotu tarczy
  \item $T_{\mathrm{p}}$ - okres precesji żyroskopu
\end{itemize}

Przekształcając równanie~(\ref{eq:zyroskop}) otrzymujemy zależność na moment bezwładności krążka metodą II ($J_{\mathrm{II}}$):
\begin{equation}
    J_{\mathrm{II}} = \frac{g}{4 \pi^2} r m T T_{\mathrm{p}}
\label{eq:moment_bezwladnosci_zyroskop}
\end{equation}

W celu uzyskania wartości $J_{\mathrm{II}}$ dokonany zostanie pomiar ramienia działania obciążenia $r$ oraz iloczynu $m T T_{\mathrm{p}}$.

Aby uzyskać wartość iloczynu $m T T_{\mathrm{p}}$, dla 9 różnych mas $m$ zostaną wykonane pomiary okresów $T$ oraz $T_{\mathrm{p}}$, w celu poznania wartości wyrażenia $\frac{1}{T T_{\mathrm{p}}}$ dla danej masy. Następnie metodą regresji liniowej uzyskane zostanie dopasowanie prostej:
\begin{equation}
    \frac{1}{T T_{\mathrm{p}}} = a m + b
\label{eq:regression_1}
\end{equation}

Wyraz wolny $b$ powinien mieć zaniedbywalną wartość. W takim przypadku równanie~(\ref{eq:regression_1}) uprości się do:
\begin{displaymath}
    \frac{1}{T T_{\mathrm{p}}} = a m
\end{displaymath}

co po przekształceniu daje następującą wartość współczynnika kierunkowego prostej $a$:
\begin{equation}
    a = \frac{1}{m T T_{\mathrm{p}}}
\label{eq:regression_2}
\end{equation}

Łącząc równości~(\ref{eq:moment_bezwladnosci_zyroskop}) i~(\ref{eq:regression_2}) otrzymujemy następującą zależność momentu bezwładności $J_{\mathrm{II}}$ od ramienia obciążenia $r$ oraz współczynnika kierunkowego prostej $a$:
\begin{equation}
    J_{\mathrm{II}}(r, a) = \frac{g}{4 \pi^2} \frac{r}{a}
\label{eq:moment_bezwladnosci_2}
\end{equation}

\section{Opracowanie danych}
\subsection{Metoda I}
\subsubsection{Średnica tarczy}
W celu obliczenia momentu bezwładności tarczy $J_{\mathrm{I}}$ pierwszą metodą, wykonane zostały pomiary średnicy tarczy $D$. Wyniki znajdują się w tabeli~(\ref{tab:srednica_tarczy}).

\begin{table}[h]
\begin{center}
\renewcommand*{\tablename}{Tabela}
\caption{\small Wyniki pomiarów średnicy tarczy $D$.\vspace{0.15cm}
    \label{tab:srednica_tarczy}}
    {\small
    \begin{tabular}{c c c}
    \hline \vspace{-0.35cm} \\
      Lp. & $D\,[\mathrm{cm}]$ & $D\,[\, \mathrm{m}]$ \\ 
      \hline \vspace{-0.2cm} \\
     1 & 25.9 & 0.259 \\
     2 & 25.3 & 0.253 \\
     3 & 24.4 & 0.244 \\
     4 & 25.7 & 0.257 \\
     5 & 25.9 & 0.259 \\
    \hline
    \end{tabular}}
    \end{center}
\end{table}


Na podstawie wyników z tabeli jesteśmy w stanie obliczyć wartość średnią:
\begin{displaymath}
    \bar{D} = 0.2544 \, \mathrm{m}
\end{displaymath}

oraz odchylenie standardowe średniej:
\begin{displaymath}
    S_{\bar{D}} \approx 2.821 \cdot 10^{-3} \, \mathrm{m}
\end{displaymath}

po uwzględnieniu współczynnika Studenta dla 5 pomiarów oraz współczynnika ufności $\alpha=0.6828$ \cite{smop_wyklad}:
\begin{displaymath}
    t_{\alpha, 5} = 1.141
\end{displaymath}

otrzymujemy niepewność standardową typu~A:
\begin{equation}
    u_{\mathrm{A}}(D) = t_{\alpha, 5} \cdot S_{\bar{D}} \approx 3.219 \cdot 10^{-3} \, \mathrm{m}
\label{eq:D_niepewnosc_A}
\end{equation}

Dokładność przyrządu (niepewność graniczna typu~B) dla pomiaru $D$ jest znana i stała, a jej wartość wynosi:
\begin{displaymath}
    \Delta D = 1 \, \mathrm{mm} = 10^{-3} \, \mathrm{m}
\end{displaymath}

na tej podstawie jesteśmy w stanie obliczyć niepewność standardową typu~B:
\begin{equation}
    u_{\mathrm{B}}(D) = \frac{\Delta D}{\sqrt{3}} \approx 5.774 \cdot 10^{-4} \, \mathrm{m}
\label{eq:D_niepewnosc_B}
\end{equation}

Na podstawie równości~(\ref{eq:D_niepewnosc_A}) oraz~(\ref{eq:D_niepewnosc_B}) jesteśmy w stanie obliczyć całkowitą niepewność standardową $D$:
\begin{displaymath}
    u(D) = \sqrt{ u_{\mathrm{A}}^2(D) + u_{\mathrm{B}}^2(D) } \approx 3.271 \cdot 10^{-3} \, \mathrm{m}
\end{displaymath}

Całkowit niepewność standardowa $u(D)$ jest bardzo bliska wartości niepewności standardowej typu~A $u_{\mathrm{A}}(D)$. Oznacza to, że niepewność standardowa typu~B $u_{\mathrm{B}}(D)$ ma jedynie nieznaczny wpływ na niepewność całkowitą.

Ostatecznie długość średnicy tarczy ma wartość:
\begin{displaymath}
    D = 0.2544(33) \, \mathrm{m}
\end{displaymath}

\subsubsection{Moment bezwładności tarczy}
Wiemy, że masa tarczy $M_{\mathrm{t}}$ ma następującą dokładną wartość:
\begin{displaymath}
    M_{\mathrm{t}} = 1.185 \, \mathrm{kg}
\end{displaymath}

Wykorzystując wzór~(\ref{eq:moment_bezwl_walca_srednica}), obliczoną wartość średnią średnicy tarczy $\bar{D}$ oraz wartość dokładną masy tarczy $M_{\mathrm{t}}$, jesteśmy w stanie obliczyć wartość oczekiwaną momentu bezwładności tarczy $J_{\mathrm{I}}$:
\begin{displaymath}
    J_{\mathrm{I}} = \frac{1}{8} M_{\mathrm{t}} D^2
\end{displaymath}
\begin{displaymath}
    \EX(J_{\mathrm{I}}) = \frac{1}{8} M_{\mathrm{t}} \bar{D}^2 \approx 9.587 \cdot  10^{-3} \, \mathrm{kg \, m^2}
\end{displaymath}

W celu obliczenia niepewności momentu bezwładności $u(J_{\mathrm{I}})$, korzystamy z metody propagacji niepewności:
\begin{equation}
    u(J_{\mathrm{I}}) = \sqrt{ [ \pdv{J_{\mathrm{I}}}{M_{\mathrm{t}}} u(M_{\mathrm{t}}) ]^2 + [ \pdv{J_{\mathrm{I}}}{D} u(D) ]^2 } \\
\label{eq:J1_niepewnosc_1}
\end{equation}

Z racji na fakt, że $M_{\mathrm{t}}$ jest wartością dokładną (a więc $u(M_{\mathrm{t}}) = 0$), wzór~(\ref{eq:J1_niepewnosc_1}) upraszcza się do postaci:
\begin{equation}
    u(J_{\mathrm{I}}) = \sqrt{ [ \pdv{J_{\mathrm{I}}}{D} u(D) ]^2 } = |\pdv{J_{\mathrm{I}}}{D} u(D)| \\
\label{eq:J1_niepewnosc_2}
\end{equation}

Pochodna cząstkowa $\pdv{J_{\mathrm{I}}}{D}$ ma postać:
\begin{equation}
    \pdv{J_{\mathrm{I}}}{D} = \frac{1}{8} M_{\mathrm{t}} \dv{D^2}{D} = \frac{1}{4} M_{\mathrm{t}} D \\
\label{eq:pochodna_dJ_dD}
\end{equation}

Łącząc równości~(\ref{eq:J1_niepewnosc_2}) i~(\ref{eq:pochodna_dJ_dD}) oraz podstawiając $D=\bar{D}$ otrzymujemy:
\begin{displaymath}
    u(J_{\mathrm{I}}) = \frac{1}{4} M_{\mathrm{t}} \bar{D}~u(D) \approx 2.465 \cdot 10^{-4}~\mathrm{kg~m^2} \\
\end{displaymath}

Ostatecznie możemy podać obliczoną wartość momentu bezwładności krążka $J_{\mathrm{I}}$ wraz z niepewnością pomiarową:
\begin{displaymath}
    J_{\mathrm{I}} = 9.59(25) \cdot 10^{-3}~\mathrm{kg~m^2} \\
\end{displaymath}

\subsection{Metoda II}
\subsubsection{Ramię działania obciążenia}
W celu obliczenia momentu bezwładności tarczy $J_{\mathrm{II}}$ drugą metodą, wykonane zostały pomiary ramienia działania obciążenia $R$. Wyniki znajdują się w tabeli~(\ref{tab:ramie_obciazenia}), a ich analiza przebiegnie identycznie jak w przypadku średnicy krążka $D$ z metody~I.

\begin{table}[h]
\begin{center}
\renewcommand*{\tablename}{Tabela}
\caption{\small Wyniki pomiarów ramienia działania obciążenia $R$.\vspace{0.15cm}
    \label{tab:ramie_obciazenia}}
    {\small
    \begin{tabular}{c c c}
    \hline \vspace{-0.35cm} \\
      Lp. & $R\,[\mathrm{cm}]$ & $R\,[\, \mathrm{m}]$ \\ 
      \hline \vspace{-0.2cm} \\
     1 & 26.4 & 0.264 \\
     2 & 27.1 & 0.271 \\
     3 & 26.9 & 0.269 \\
     4 & 26.6 & 0.266 \\
     5 & 26.6 & 0.266 \\
    \hline
    \end{tabular}}
    \end{center}
\end{table}


Wartość średnia:
\begin{displaymath}
    \bar{R} = 0.2672 \, \mathrm{m}
\end{displaymath}

oraz odchylenie standardowe średniej:
\begin{displaymath}
    S_{\bar{R}} \approx 1.241 \cdot 10^{-3} \, \mathrm{m}
\end{displaymath}

niepewność standardową typu~A po uwzględnieniu współczynnika Studenta:
\begin{equation}
    u_{\mathrm{A}}(R) = t_{\alpha, 5} \cdot S_{\bar{R}} \approx 1.416 \cdot 10^{-3} \, \mathrm{m}
\label{eq:R_niepewnosc_A}
\end{equation}

Dokładność przyrządu (niepewność graniczna typu~B) dla pomiarów $R$ jest identyczna jak w przypadku pomiaru $D$, a więc także niepewność standardowa typu~B będzie miała identyczną wartość:
\begin{equation}
    u_{\mathrm{B}}(R) = u_{\mathrm{B}}(D) \approx 5.774 \cdot 10^{-4} \, \mathrm{m}
\label{eq:R_niepewnosc_B}
\end{equation}

Na podstawie równości~(\ref{eq:R_niepewnosc_A}) oraz~(\ref{eq:R_niepewnosc_B}) jesteśmy w stanie obliczyć całkowitą niepewność standardową $R$:
\begin{displaymath}
    u(R) = \sqrt{ u_{\mathrm{A}}^2(R) + u_{\mathrm{B}}^2(R) } \approx 1.529 \cdot 10^{-3} \, \mathrm{m} \\
\end{displaymath}

W przypadku pomiarów $R$ niepewność typu~A jest także znacznie większa, niż typu~B, jednak różnica nie jest aż tak wielka jak w wypadku pomiarów $D$. Niepewność typu~B $u_{\mathrm{B}}(D)$ zmienia w sposób istotny niepewność całkowitą.

Ostatecznie długość ramienia działania obciążenia ma wartość:
\begin{displaymath}
    R = 0.2672(16) \, \mathrm{m} \\
\end{displaymath}

\subsubsection{Pomiary okresów dla danych mas}

Dla każdej z 9 różnych mas($m_1, m_2, \dots, m_9$) wykonano po 10 pomiarów trzech odcinków czasu: $T_1, T_2, T_3$. Z tych odcinków można obliczyć okres obrotu $T$:
\begin{displaymath}
    T = \frac{T_1 + T_2}{2} \\
\end{displaymath}

oraz okres precesji $T_{\mathrm{p}}$:
\begin{displaymath}
    T_{\mathrm{p}} = 2 T_3 \\
\end{displaymath}

Przykładowe pomiary odcinków czasu (dla $m_1 = 10 \, \mathrm{g}$) wraz obliczoną odpowiadającą wartością $p = \frac{1}{T T_{\mathrm{p}}}$ znajdują się w tabeli (\ref{tab:pomiary_czas_10g}).

\begin{table}[h]
\begin{center}
\renewcommand*{\tablename}{Tabela}
\caption{\small Przykładowe wyniki pomiarów odcinków czasu dla $m_1 = 10 \, \mathrm{g}$ wraz z obliczonym wyrażeniem~$\frac{1}{T T_{\mathrm{p}}}$.\vspace{0.15cm}
    \label{tab:pomiary_czas_10g}}
    {\small
    \begin{tabular}{c c c c c}
    \hline \vspace{-0.35cm} \\
      Lp. & $T_1\,[\mathrm{s}]$ & $T_2\,[\mathrm{s}]$ & $T_3\,[\mathrm{s}]$ & $10^2 \frac{1}{T T_{\mathrm{p}}}\,[\mathrm{s}^{-2}]$ \\ 
      \hline \vspace{-0.2cm} \\
     1 & 0.2506 & 0.2438 & 21.06 & 9.604 \\
     2 & 0.3945 & 0.3752 & 23.46 & 5.538 \\
     3 & 0.1990 & 0.1950 & 23.75 & 10.69 \\
     4 & 0.3936 & 0.3818 & 25.29 & 5.099 \\
     5 & 0.2845 & 0.2735 & 19.53 & 9.176 \\
     6 & 1.1531 & 1.1140 & 13.09 & 3.37 \\
     7 & 1.2096 & 1.1748 & 13.06 & 3.211 \\
     8 & 0.4788 & 0.4647 & 19.56 & 5.419 \\
     9 & 0.1596 & 0.1528 & 22.29 & 14.36 \\
     10 & 0.3820 & 0.3673 & 25.04 & 5.328 \\
    \hline
    \end{tabular}}
    \end{center}
\end{table}


Następnie dla każdego $m_i$ liczymy wartość średnią wyrażenia $p = \frac{1}{T T_{\mathrm{p}}}$ oraz odchylenie standardowe średniej $S_p$.

Uwzględniamy także współczynnik Studenta dla 10 pomiarów oraz współczynnika ufności $\alpha=0.6828$ \cite{smop_wyklad}:
\begin{displaymath}
    t_{\alpha, 10} = 1.059
\end{displaymath}

Co daje postać niepewności standardowej typu~A:
\begin{displaymath}
    u_{\mathrm{A}}(p) = t_{\alpha, 10} \cdot S_{\bar{p}}
\end{displaymath}

Obliczone wartości wyrażenia $p = \frac{1}{T T_{\mathrm{p}}}$ wraz z odchyleniem standardowym $S_p$ oraz niepewnością standardową typu~A $u_{\mathrm{A}}(p)$ dla każdego $m_i$ znajdują się w tabeli (\ref{tab:regresja_dane}).

\begin{table}[h]
\begin{center}
\renewcommand*{\tablename}{Tabela}
\caption{\small Obliczone wartości wyrażenia $p = \frac{1}{T T_{\mathrm{p}}}$ wraz z odchyleniem standardowym $S_p$ oraz niepewnością standardową typu~A $u_{\mathrm{A}}(p)$ dla każdego $m_i$.\vspace{0.15cm}
    \label{tab:regresja_dane}}
    {\small
    \begin{tabular}{c c c c c}
    \hline \vspace{-0.35cm} \\
      Lp. & $m_i \, [\mathrm{g}]$ & $10p \, [\mathrm{s}^{-2}]$ & $10 S_p \, [\mathrm{s}^{-2}]$ & $10 u_{\mathrm{A}}(p) \, [\mathrm{s}^{-2}]$ \\ 
      \hline \vspace{-0.2cm} \\
     1 & 10 & 0.7179 & 0.1143 & 0.121 \\
     2 & 20 & 2.21 & 0.1995 & 0.2113 \\
     3 & 30 & 3.216 & 0.4075 & 0.4316 \\
     4 & 40 & 4.448 & 0.3194 & 0.3382 \\
     5 & 60 & 4.771 & 0.528 & 0.5592 \\
     6 & 70 & 5.43 & 0.4927 & 0.5218 \\
     7 & 80 & 6.044 & 1.174 & 1.2436 \\
     8 & 90 & 6.701 & 0.3296 & 0.3491 \\
     9 & 100 & 8.272 & 0.4195 & 0.4442 \\
    \hline
    \end{tabular}}
    \end{center}
\end{table}


Dokładność przyrządu (niepewność graniczna typu~B) dla pomiaru $T_1$ i $T_2$ jest znana i stała, a jej wartość wynosi:
\begin{displaymath}
    \Delta T_1 = \Delta T_2 = 10^{-4} \, \mathrm{s}
\end{displaymath}
\begin{displaymath}
    \Delta T_3 = 10^{-2} \, \mathrm{s}
\end{displaymath}

Jak widać w tabeli (\ref{tab:regresja_dane}) obliczone niepewności standardowe typu~A $u_{\mathrm{A}}(p)$ są rzędu $10 \%$ wartości. Dla porównania niepewności graniczne typu~B są rzędu $10^{-4} - 10^{-3}$ ($T_1$ i $T_2$) oraz ok. $10^{-3}$ dla $T_3$. Ze względu na tak dużą dysproporcję, niepewności pomiarowe typu~B mogą zostać pominięte w obliczeniach, a więc:
\begin{displaymath}
    u(p) = u_{\mathrm{A}}(p)
\end{displaymath}

\subsubsection{Regresja liniowa}
Dane z tabeli (\ref{tab:regresja_dane}) przedstawione są na wykresie (Rysunek \ref{img:regresja_wykres}). Następnie wykonana została regresja liniowa przy pomocy programu \textit{OriginLab} \footnote{\url{https://www.originlab.com/}}. Użyta została regresja ważona w wariancie \textit{Instrumental}, gdzie waga danego punktu pomiarowego jest dana jako \cite{origin_fitting_weighting}:
\begin{displaymath}
    w_i = \frac{1}{\sigma_i^2}
\end{displaymath}

gdzie $\sigma_i$ jest wartością $u_{\mathrm{A}}(p)$ z tabeli (\ref{tab:regresja_dane}) dla danego punktu.

\begin{figure}[h]
    \centering

    \begin{tikzpicture}
        \begin{axis}[
            width=8.4cm, height=7cm,
            xlabel={$m \, [\mathrm{g}]$},
            ylabel={$p = \frac{1}{T T_p} \, [\mathrm{s}^{-2}]$},
            legend pos=north west,
            grid=major
        ]

        \addlegendimage{only marks, mark=*, blue, mark options={fill=red}, error bars/.cd, y dir=both, y explicit}
        \addlegendentry{Punkty pomiarowe}
        
        \addlegendimage{no markers, red, thick}
        \addlegendentry{Dopasowanie regresji liniowej}
        
        \def\slope{0.00795279}
        \def\intercept{0.01695}
        
        \addplot [red, thick, domain=10:100, forget plot] {\slope*x + \intercept};
        \addlegendentry{Dopasowanie regresji liniowej}
    
        \addplot+[
            only marks,
            mark=*,
            blue,
            mark options={fill=red},
            forget plot,
            error bars/.cd,
            y dir=both,
            y explicit,
        ] table [x=x, y=y, y error=err] {
            x   y       err
            10	0.07179	0.0121
            20	0.22102	0.02113
            30	0.32161	0.04316
            40	0.44478	0.03382
            60	0.47714	0.05592
            70	0.54298	0.05218
            80	0.60435	0.12436
            90	0.67012	0.03491
            100	0.82715	0.04442
        };
        \addlegendentry{Punkty pomiarowe}
    
        \end{axis}
    \end{tikzpicture}

    \caption{Masy $m_i$ wraz z odpowiadającymi im średnimi wartościami wyrażenia $p = \frac{1}{T T_p}$, policzonymi niepewnościami standardowymi $u(p)$ oraz dopasowaniem krzywej $y = ax + b$ uzyskanej dzięki regresji liniowej.}
    \label{img:regresja_wykres}
\end{figure}
    

W wyniku regresji otrzymane zostały następujące wartości parametrów prostej:
\begin{displaymath}
\begin{split}
    a &= (7.95279 \pm 0.62226) \, \mathrm{kg^{-1} \, s^{-2}} \\
    &\approx 7.95(63) \, \mathrm{kg^{-1} \, s^{-2}}
\end{split}
\end{displaymath}
\begin{displaymath}
    b = (1.695 \pm 2.391) \cdot 10^{-2} \, \mathrm{s^{-2}}
\end{displaymath}

Krzywa, wraz z danymi pomiarowymi znajduje się na wykresie (Rysunek \ref{img:regresja_wykres}).

Miara jakości dopasowania \textit{Adjusted R-Square} \cite{origin_adj_r_square} ma następującą wartość, co uznawane jest za \enquote{dobre} dopasowanie:
\begin{displaymath}
    R^2_{\mathrm{adj}} = 0.95304
\end{displaymath}

Zgodnie z oczekiwaniami, zerowa wartość wyrazu wolnego $b$ mieści się w granicach błędu statystycznego dopasowania, a więc zostanie on pominięty w dalszych rozważaniach (uproszczenie zależności (\ref{eq:regression_1}) do (\ref{eq:regression_2})).

Błąd standardowy współczynnika kierunkowego $a$ jest rzędu $10\%$ wartości, a więc ponownie błąd standardowy typu~B zostanie pominięty w dalszych rozważaniach, a wartość błędu $S_a$ uzyskana z regresji będzie pełniła rolę całkowitej niepewności standardowej $u(a)$:
\begin{displaymath}
    u(A) = S_a = 0.62226 \, \mathrm{kg^{-1} \, s^{-2}}
\end{displaymath}

\subsubsection{Moment bezwładności tarczy}
Korzystając ze wzoru (\ref{eq:moment_bezwladnosci_2}), obliczonej długości ramienia $R$ oraz współczynnika kierunkowego $a$ otrzymujemy wzór na moment bezwładności tarczy metodą drugą $J_{\mathrm{II}}$:
\begin{equation}
    J_{\mathrm{II}} = J_{\mathrm{II}}(R, a) = \frac{g R}{4 \pi^2 a}
\label{eq:J2_wzor}
\end{equation}

Przyspieszenie ziemskie na terenie Krakowa ma wartość \cite{przyspieszenie_ziemskie_krakow}:
\begin{displaymath}
    g_{\mathrm{K}} = 9.8100(5) \, \mathrm{\frac{m}{s^2}}
\end{displaymath}

W celu obliczenia wartości oczekiwanej $\EX(J_{\mathrm{II}})$, podstawiamy do wyrażenia (\ref{eq:J2_wzor}) średnią długość ramienia $\bar{R}$, wartość oczekiwaną współczynnika kierunkowego $\EX(a)$ oraz wartość oczekiwaną przyspieszenia ziemskiego $\EX(g_{\mathrm{K}})$
\begin{displaymath}
\begin{split}
    \EX(J_{\mathrm{II}})
    &= \frac{\EX(g_{\mathrm{K}}) \bar{R}}{4 \pi^2 \EX(a)} \\
    &= \frac{9.8100 \, \mathrm{m \, s^{-2}} \cdot 0.2672 \, m}
    {4 \pi^2 \cdot 7.95279 \, \mathrm{kg^{-1} \, s^{-2}}} \\
    &\approx 8.34884 \cdot 10^{-3} \, \mathrm{kg \, m^2}
\end{split}
\end{displaymath}

W celu obliczenia niepewności momentu bezwładności $u(J_{\mathrm{II}})$, korzystamy z metody propagacji niepewności:
\begin{displaymath}
\begin{split}
    u(J_{\mathrm{II}}) &= \\
    =& \sqrt{
        [ \pdv{J_{\mathrm{II}}}{g} u(g) ]^2
        + [ \pdv{J_{\mathrm{II}}}{R} u(R) ]^2
        + [ \pdv{J_{\mathrm{II}}}{a} u(a) ]^2
    }
\end{split}
\end{displaymath}

Poszczególne pochodne cząstkowe mają postać:
\begin{displaymath}
    \pdv{J_{\mathrm{II}}}{g} = \frac{R}{4 \pi^2 a}
\end{displaymath}

\begin{displaymath}
    \pdv{J_{\mathrm{II}}}{R} = \frac{g}{4 \pi^2 a}
\end{displaymath}

\begin{displaymath}
    \pdv{J_{\mathrm{II}}}{a} = -\frac{g R}{4 \pi^2 a^2}
\end{displaymath}

a więc finalnie wzór na niepewność momentu bezwładności $u(J_{\mathrm{II}})$ będzie miał podstać:
\begin{equation}
\begin{split}
    u(J_{\mathrm{II}}) &= \\
    =& \sqrt{
        ( \frac{g R}{4 \pi^2 a} )^2 [
            ( \frac{u(g)}{g} )^2
            + ( \frac{u(R)}{R} )^2
            + ( \frac{u(a)}{a} )^2
        ]
    } \\
    =& \frac{g R}{4 \pi^2 a} \sqrt{
        ( \frac{u(g)}{g} )^2
        + ( \frac{u(R)}{R} )^2
        + ( \frac{u(a)}{a} )^2
    }
\end{split}
\end{equation}

Poszczególne składniki sumy pod pierwiastkiem mają w przybliżeniu wartości:
\begin{displaymath}
    ( \frac{u(g)}{g} )^2 \approx 2.6 \cdot 10^{-9}
\end{displaymath}

\begin{displaymath}
    ( \frac{u(R)}{R} )^2 \approx 3.28 \cdot 10^{-5}
\end{displaymath}

\begin{displaymath}
    ( \frac{u(a)}{a} )^2 \approx 6.12 \cdot 10^{-3}
\end{displaymath}

Widzimy więc, że to niepewność współczynnika kierunkowego prostej $u(a)$ odgrywa zdecydowanie decydującą rolę przy liczeniu niepewności momentu bezwładności $u(J_{\mathrm{II}})$.

Finalnie, podstawiając wszystkie obliczone wartości uzyskujemy:
\begin{displaymath}
    u(J_{\mathrm{II}}) \approx 4.115 \cdot 10^{-3} \, \mathrm{kg \, m^2}
\end{displaymath}

Co daje następującą wartość $J_{\mathrm{II}}$ wraz z niepewnością:
\begin{displaymath}
    J_{\mathrm{II}} = (8.3 \pm 4.2) \cdot 10^{-3} \, \mathrm{kg \, m^2}
\end{displaymath}

\section{Podsumowanie}
W wyniku obliczeń obie metody dały następujące wyniki:
\begin{displaymath}
    J_{\mathrm{I}} = 9.59(25) \cdot 10^{-3}~\mathrm{kg~m^2} \\
\end{displaymath}

\begin{displaymath}
    J_{\mathrm{II}} = (8.3 \pm 4.2) \cdot 10^{-3} \, \mathrm{kg \, m^2}
\end{displaymath}

Funkcje gęstości prawdopodobieństwa dla rozkładów wartości momentów bezwładności dla obu metod znajdują się na wykresach (Rysunek \ref{img:gestosc_prawdop_porownanie_szeroko}) - szeroka domena, oraz (Rysunek \ref{img:gestosc_prawdop_porownanie_wasko}) - wąska domena. Zaznaczone zostały również wartości oczekiwane $\EX(J_{\mathrm{I}})$ oraz $\EX(J_{\mathrm{II}})$.

\begin{figure}[h]
    \centering
    \begin{tikzpicture}
        \begin{axis}[
            width=8.2cm, height=7cm,
            xlabel={$10^{-3} J \, [kg \, m^2]$},
            ylabel={Gęstość prawdopodobieństwa},
            grid=major,
            scaled x ticks=false,
            x tick label style={/pgf/number format/fixed, /pgf/number format/precision=4},
            legend style={at={(0.02,0.98)}, anchor=north west}
        ]

        % --- Metoda I ---
        \def\muOne{9.59}
        \def\sigmaOne{0.25}

        % --- Metoda II ---
        \def\muTwo{8.3}
        \def\sigmaTwo{4.2}

        \def\medium{(1 * \muOne + 1 * \muTwo) / 2}
        \def\range{2 * \sigmaTwo}

        % --- Metoda I ---
        \addplot [
            domain=\medium - \range : \medium + \range,
            samples=1000,
            color=red,
            thick,
            smooth
        ]
        {1/(\sigmaOne*sqrt(2*pi)) * exp(-((x-\muOne)^2)/(2*\sigmaOne^2))};
        \addlegendentry{Metoda I}

        % --- Metoda II ---
        \addplot [
            domain=\medium - \range : \medium + \range,
            samples=1000,
            color=blue,
            thick,
            smooth
        ]
        {1/(\sigmaTwo*sqrt(2*pi)) * exp(-((x-\muTwo)^2)/(2*\sigmaTwo^2))};
        \addlegendentry{Metoda II}

        % --- Punkty --- %
        \pgfmathsetmacro{\peakOne}{1/(\sigmaOne*sqrt(2*pi))}
        \draw [dashed, black!70] (axis cs:\muOne, 0) -- (axis cs:\muOne, \peakOne);
        \addplot[mark=*, mark size=1.5pt, red] coordinates {(\muOne, \peakOne)};
        \addplot[mark=*, mark size=1.5pt, red] coordinates {(\muOne, 0)};
        
        \pgfmathsetmacro{\peakTwo}{1/(\sigmaTwo*sqrt(2*pi))}
        \draw [dashed, black!70] (axis cs:\muTwo, 0) -- (axis cs:\muTwo, \peakTwo);
        \addplot[mark=*, mark size=1.5pt, blue] coordinates {(\muTwo, \peakTwo)};
        \addplot[mark=*, mark size=1.5pt, blue] coordinates {(\muTwo, 0)};

        \end{axis}
    \end{tikzpicture}
    \caption{Porównanie gęstości prawdopodobieństwa uzyskanych momentów bezwładności dla obu metod wraz z zaznaczonymi wartościami oczekiwanymi (szeroka domena).}
    \label{img:gestosc_prawdop_porownanie_szeroko}
\end{figure}


\begin{figure}[h]
    \centering
    \begin{tikzpicture}
        \begin{axis}[
            width=8.2cm, height=7cm,
            xlabel={$10^{-3} J \, [kg \, m^2]$},
            ylabel={Gęstość prawdopodobieństwa},
            grid=major,
            scaled x ticks=false,
            x tick label style={/pgf/number format/fixed, /pgf/number format/precision=4},
            legend style={at={(0.02,0.98)}, anchor=north west}
        ]

        % --- Metoda I ---
        \def\muOne{9.59}
        \def\sigmaOne{0.25}

        % --- Metoda II ---
        \def\muTwo{8.3}
        \def\sigmaTwo{4.2}

        \def\medium{(2 * \muOne + 1 * \muTwo) / 3}
        \def\range{4 * \sigmaOne}

        % --- Metoda I ---
        \addplot [
            domain=\medium - \range : \medium + \range,
            samples=1000,
            color=red,
            thick,
            smooth
        ]
        {1/(\sigmaOne*sqrt(2*pi)) * exp(-((x-\muOne)^2)/(2*\sigmaOne^2))};
        \addlegendentry{Metoda I}

        % --- Metoda II ---
        \addplot [
            domain=\medium - \range : \medium + \range,
            samples=1000,
            color=blue,
            thick,
            smooth
        ]
        {1/(\sigmaTwo*sqrt(2*pi)) * exp(-((x-\muTwo)^2)/(2*\sigmaTwo^2))};
        \addlegendentry{Metoda II}

        % --- Punkty --- %
        \pgfmathsetmacro{\peakOne}{1/(\sigmaOne*sqrt(2*pi))}
        \draw [dashed, black!70] (axis cs:\muOne, 0) -- (axis cs:\muOne, \peakOne);
        \addplot[mark=*, mark size=1.5pt, red] coordinates {(\muOne, \peakOne)};
        \addplot[mark=*, mark size=1.5pt, red] coordinates {(\muOne, 0)};
        
        \pgfmathsetmacro{\peakTwo}{1/(\sigmaTwo*sqrt(2*pi))}
        \draw [dashed, black!70] (axis cs:\muTwo, 0) -- (axis cs:\muTwo, \peakTwo);
        \addplot[mark=*, mark size=1.5pt, blue] coordinates {(\muTwo, \peakTwo)};
        \addplot[mark=*, mark size=1.5pt, blue] coordinates {(\muTwo, 0)};

        \end{axis}
    \end{tikzpicture}
    \caption{Porównanie gęstości prawdopodobieństwa uzyskanych momentów bezwładności dla obu metod wraz z zaznaczonymi wartościami oczekiwanymi (wąska domena).}
    \label{img:gestosc_prawdop_porownanie_wasko}
\end{figure}


Na szczególną uwagę zasługuje ogromna niepewność wyniku w metodzie II. Błąd względny wynosi:
\begin{displaymath}
    \frac{ u(J_{\mathrm{II}}) }{ \EX(J_{\mathrm{II}}) } = 49.3\%
\end{displaymath}

Przez co wartość maksymalna przedziału ufności dla $\alpha = 68.28\%$ jest niemalże trzykrotnością wartości minimalnej:
\begin{displaymath}
    \frac{ \EX(J_{\mathrm{II}}) + u(J_{\mathrm{II}}) }{ \EX(J_{\mathrm{II}}) - u(J_{\mathrm{II}}) }
    \approx 2.94
\end{displaymath}

Dla porównania w metodzie I ograniczona niepewność daje znacznie większe pole do analizy obliczonej wartości:
\begin{displaymath}
    \frac{ u(J_{\mathrm{I}}) }{ \EX(J_{\mathrm{I}}) } = 2.57\%
\end{displaymath}

\begin{displaymath}
    \frac{ \EX(J_{\mathrm{I}}) + u(J_{\mathrm{I}}) }{ \EX(J_{\mathrm{I}}) - u(J_{\mathrm{I}}) }
    \approx 1.05
\end{displaymath}

\begin{thebibliography}{7}
\bibitem{pracownia_fizyczna_m8} I~Pracownia Fizyczna, \emph{Badanie ruchu precesyjnego żyroskopu (M8)}, Instytut Fizyki~UJ (\url{https://1pf.if.uj.edu.pl/materialy-do-cwiczen/spis-cwiczen}).
\bibitem{eFizyka_bryla_sztywna} prof.~dr~hab.~inż. \mbox{Zbigniew} \mbox{Kąkol}, \mbox{dr~Jan} \mbox{Żukrowski}, \emph{e-Fizyka}, rozdział \textbf{11.3} \emph{Dynamika bryły sztywnej} (\url{https://home.agh.edu.pl/~kakol/efizyka/w11/main11c.html}).
\bibitem{smop_wyklad} dr~hab. \mbox{Katarzyna} \mbox{Gajos}, Wykład \emph{Statystyczne Metody Opracowania Pomiarów I Dla studentów fizyki}.
\bibitem{origin_fitting_weighting} \emph{Origin Help}, rozdział \textbf{15.3.6.5} \emph{Fitting with Errors and Weighting}, OriginLab Corporation (\url{https://www.originlab.com/doc/Origin-Help/FIt-with-Err-Weight}).
\bibitem{origin_adj_r_square} \emph{Origin Help}, rozdział \textbf{15.4} \emph{Interpreting Regression Results}, sekcja \textbf{2.5} \emph{Adj. R-Square}, OriginLab Corporation (\url{https://www.originlab.com/doc/Origin-Help/Interpret-Regression-Result#Adj._R-Square}).
\bibitem{przyspieszenie_ziemskie_krakow} \emph{Dziennik normalizacji i miar nr 13/81 r.}, Polski Komitet Normalizacji, Miar i Jakości (\url{https://www.gum.gov.pl/download/2/1854/dzienniknormalizacjiimiar198113.pdf}).

\end{thebibliography}

\end{document}
